\notechapter{N-20221217}{Function Exercises: Domains, Ranges, Dirichlet Function, Parameters, and Functional Equations}

\section{Structure of the exercise sheet}

The note is a function exercise sheet, not a single theorem.  It contains domain problems, range problems, a Dirichlet-function example, parameter problems, zero-counting by graphs, and functional equations.  The repaired version keeps the exercise logic: each problem is a training example for a specific method.

\section{Domain problems}

For a function domain, the first rule is to list all legality conditions.  Square roots require nonnegative radicands, logarithms require positive arguments and valid bases, denominators cannot be zero, and iterated functions require every intermediate input to lie in the next domain.

For example,
\[
  y=\frac{\sqrt{x(3-x)}}{\lg (x-3)^2}
\]
requires
\[
  x(3-x)\ge0,
  \qquad
  (x-3)^2>0,
  \qquad
  \lg (x-3)^2\ne0.
\]
The first gives $0\le x\le3$.  The second excludes $x=3$.  The third excludes $(x-3)^2=1$, so $x=2$ or $4$, but only $2$ lies in the current interval.  Hence
\[
  D=[0,3)\setminus\{2\}.
\]
The important step is not the answer; it is the intersection of all restrictions.

For an expression such as
\[
  \sqrt{a^x-kb^x}
\]
in a denominator, the condition is strict:
\[
  a^x-kb^x>0.
\]
Since $b^x>0$, divide by $b^x$:
\[
  \left(\frac ab\right)^x>k.
\]
Now the inequality depends on whether $a/b>1$ or $0<a/b<1$.  This is a recurring trap in exponential and logarithmic domain problems.

\section{Floor functions}

The exercise sheet also uses the Gauss or floor function $[x]$.  The correct method is interval splitting:
\[
  [x]=m \quad\Longleftrightarrow\quad m\le x<m+1.
\]
A floor-function domain problem is therefore a finite or countable case split.  One should not treat $[x]$ as a smooth algebraic expression.

\section{Range problems}

For
\[
  y=\sqrt{x-4}+\sqrt{15-3x},
\]
the domain is $4\le x\le5$.  Put $u=x-4$, so $0\le u\le1$ and
\[
  y=\sqrt u+\sqrt{3(1-u)}.
\]
By Cauchy--Schwarz,
\[
  y^2\le (1+3)(u+1-u)=4,
\]
so $y\le2$.  At the endpoints the smaller value is $1$, and the maximum is achieved when the Cauchy equality condition holds.  Thus the range is $[1,2]$.

For
\[
  y=x+\sqrt{1-2x},
\]
let $t=\sqrt{1-2x}\ge0$.  Then
\[
  x=\frac{1-t^2}{2},\qquad
  y=-\frac12(t-1)^2+1.
\]
Hence the range is
\[
  (-\infty,1].
\]
This substitution method is the natural way to remove a radical while preserving the domain.

\section{Dirichlet function}

The Dirichlet function is
\[
  D(x)=\begin{cases}
  1,&x\in\mathbb Q,\\
  0,&x\in\mathbb R\setminus\mathbb Q.
  \end{cases}
\]
It is bounded, because $0\le D(x)\le1$.  It is not monotone on any interval, because every interval contains both rational and irrational numbers.  It has every nonzero rational period $T$, since
\[
  x\in\mathbb Q \Longleftrightarrow x+T\in\mathbb Q
\]
for $T\in\mathbb Q$.  It has no irrational period, because adding an irrational changes rationality.

This exercise is included because it teaches that functions can be extremely discontinuous while still having simple algebraic properties.

\section{Graph transformations and zero counting}

The problem involving $f$ and
\[
  g(x)=b-f(2-x)
\]
asks for the values of $b$ for which $f(x)-g(x)$ has four zeros.  The transformation $x\mapsto2-x$ reflects the graph across the vertical line $x=1$; then $b-$ reflects vertically and shifts.  The number of zeros of $f-g$ is the number of intersections of two transformed graphs.

The method is graphical but not vague.  The number of intersections changes only when a moving graph becomes tangent, passes through a vertex/corner, or crosses an endpoint.  Thus parameter intervals are found by checking these critical positions.

\section{Logarithmic parameter equation}

For an equation such as
\[
  \ln(ax)=2\ln(x+1),
\]
first impose the domain:
\[
  ax>0,
  \qquad x+1>0.
\]
Then combine logarithms:
\[
  \ln(ax)=\ln((x+1)^2)
\]
so
\[
  ax=(x+1)^2.
\]
This becomes a quadratic in $x$.  Having exactly one real solution means a discriminant or tangency condition, but the domain must still be checked.  The sheet is training this two-step habit: algebraic transformation plus domain verification.

\section{Functional equations with period-like recurrence}

Suppose $f$ is defined on $[0,1)$ by
\[
  f(x)=2x-x^2
\]
or another explicit expression, and for all real $x$ satisfies
\[
  f(x)+f(x+1)=1.
\]
Then values outside $[0,1)$ are reduced by repeatedly using
\[
  f(x+1)=1-f(x).
\]
If $a=\log_2 3$, to compute $f(a)+f(2a)+f(3a)$, first determine the integer intervals containing $a,2a,3a$, reduce their fractional parts into $[0,1)$, and then apply the recurrence each time.

This is the same idea as reducing modulo a period, except that shifting by $1$ may flip the value rather than preserve it.

\section{A reciprocal functional equation}

The sheet also has a functional equation of the form
\[
  f(x)=f(1)x+f(2)\frac1x-1
\]
for nonzero $x$.  On $(0,+\infty)$, this becomes a problem about minimizing
\[
  Ax+\frac Bx-1.
\]
The AM--GM inequality gives
\[
  Ax+\frac Bx\ge2\sqrt{AB}
\]
when $A,B>0$.  The equality condition determines the point where the minimum occurs.  This is a typical example where a functional equation reduces to a standard inequality.

\section{Quadratic iteration problem}

For a quadratic
\[
  f(x)=x^2+2ax+8,
\]
conditions involving both $f(x)\le0$ and $f(f(x))\le8$ should be treated by introducing the image variable $u=f(x)$.  Then the second condition becomes $f(u)\le8$.  The problem is about how the range of the first inequality is mapped by the quadratic again.  Completing the square and tracking intervals is safer than expanding $f(f(x))$ blindly.

\section{Max and absolute value}

The exercise about
\[
  M(x)=\max\{f(x),g(x)\}
\]
uses the identity
\[
  \max\{u,v\}=\frac{u+v+|u-v|}{2}.
\]
Thus maximum and minimum operations can be encoded using absolute value.  This is useful when deciding whether a function built by max/min remains elementary.


\section{Parameter problems and critical values}

When a parameter changes the number of solutions, the change usually happens at a critical value: a tangent contact, an endpoint contact, or a collision of two roots.  This is why discriminants and graph tangencies appear so often.

For a quadratic equation depending on a parameter, the condition for exactly one real root is usually
\[
  \Delta=0.
\]
For a graphical equation, the same condition appears geometrically as tangency.  The algebraic and geometric languages are two versions of the same idea.

\section{Functional equations as reduction rules}

A functional equation such as
\[
  f(x+1)=1-f(x)
\]
should be used as a reduction rule.  It tells us how to move an input back to a fundamental interval.  After two shifts,
\[
  f(x+2)=1-f(x+1)=f(x),
\]
so the function becomes periodic with period \(2\).  This kind of hidden periodicity is easy to miss if one treats the equation only algebraically.

\section{Expanded synthesis and advanced context}

The page is not just a set of answers.  It teaches a toolbox: domains are legality intersections; ranges use substitution and inequalities; floor functions require interval splitting; Dirichlet's function uses density of rationals and irrationals; parameter problems are tangency/discriminant problems; functional equations reduce values to a fundamental interval; max and min can be represented by absolute values.  The common theme is that a function is not merely a formula but a domain, graph, transformation rule, and set of constraints.


\paragraph{Expanded synthesis.}
The elementary geometric content of this chapter should be read through transformations and invariants. The earlier sections (`Structure of the exercise sheet`, `Domain problems`, `Floor functions`, `Range problems`, `Dirichlet function`) may use coordinates, angles, determinants, parametrizations, or integral formulas, but the organizing question is: which quantities survive the transformations naturally associated with the problem? Euclidean geometry preserves distances and angles; affine geometry preserves lines and ratios on a line; projective geometry preserves incidence and cross ratio; differential geometry preserves coordinate-independent objects such as forms, metrics, and orientations.

This perspective separates different layers of a problem. A dot product belongs to metric geometry. A determinant belongs to oriented volume. A cross ratio belongs to projective geometry. A line or surface integral of a differential form belongs to oriented topology. A scalar area integral belongs to the metric-induced measure. Confusing these layers is the source of many errors, especially in vector calculus where $dS$, $F\cdot n\,dS$, $dx\wedge dy$, and boundary orientation look similar but encode different structures.

When returning to the original examples, identify the natural object before computing. If the problem is about concurrence or collinearity, projective or affine tools may be natural. If it is about distance, angle, or orthogonality, metric tools are essential. If it is about flux or circulation, differential forms and orientation explain the signs. The advanced viewpoint makes the elementary solution more disciplined rather than more abstract for its own sake.
